\begin{textsample}{POS Dim 3 – human – Score 6.00 – sp/sp\_inf\_e\_ef\_intro\_2\_p7.txt}  \label{ex:f3_pos_010}
Em 1996, tem início o processo de municipalização do Ensino Fundamental, \textbf{etapa} até então inteiramente sob a responsabilidade do \textbf{Estado}.
Naquele ano, 46 municípios paulistas iniciaram o processo de municipalização do ensino, assumindo — por meio de parceria Estado-Município — as \textbf{primeiras} séries do 1º grau, ampliando as \textbf{etapas} até então atendidas, \textbf{visto} que parte significativa dos municípios oferecia apenas a Educação Infantil ( Pré-Escola ) em suas redes.
Vale destacar que as creches, até os anos 2000, eram vinculadas, nos munícipios, à instância do bem-estar social.
A partir desse período, a Educação Infantil passou a integrar a rede de escolas das Secretarias Municipais de Educação, com a ampliação dos segmentos creche e pré-escola.

% matched lemmas: estado, etapa, primeiro, ver
\end{textsample}
